\chapter{Methodology}
\label{methodology}

This chapter presents a collection of probabilistic models used to describe football match outcomes.
Rather than introducing each model independently, we adopt a unified approach in which both model families are first defined through general formulations.
Specific models are then obtained by imposing structural constraints and simplifying assumptions on these general representations.

We begin with the Bradley–Terry family of models, which describe match outcomes in terms of latent team skill parameters.
A general formulation is introduced that encompasses home effects and the possibility of draws, from which simpler variants naturally emerge by restricting the corresponding parameters.
This formulation provides a common reference point for interpretation, comparison, and validation across the different Bradley–Terry specifications.

We then consider Poisson-based models, which explicitly model the number of goals scored by each team.
The general Poisson formulation incorporates offensive and defensive effects, home effects, and potential dependence between teams’ scoring processes.
As with the Bradley–Terry case, the specific Poisson models used in this work are obtained through systematic simplifications of the general formulation, resulting in a structured hierarchy of models with increasing complexity.

All models are intentionally kept simple, both in structure and in the number of parameters.
This design choice reflects our goal of validating the models using Simulation-Based Calibration and fitting them to real data from the Brazilian First Division.
The predictive performance of the models is compared, and the best-performing specifications are subsequently extended to player-level versions, where the dimensionality of the parameter space increases substantially.
Maintaining simple club-level formulations is therefore essential to control the computational cost of these extensions while preserving interpretability.

The chapter also describes how the club-level formulations are adapted to player-level models, outlining the structural modifications required to replace team-level to player-level parameters.
The inference procedure is then detailed, covering the prior distributions, the data inputs at each level, and the MCMC sampling configuration.
Finally, the chapter concludes with the evaluation methodology used to assess the predictive performance of the models on real data, including the prediction scheme, baseline predictors, and proper scoring rules.

\section{Bradley-Terry model}
\label{sec:BTmodel}

Bradley–Terry models \cite{bradley_terry_1952} provide a natural starting point for modelling football match outcomes, as they directly relate the probability of winning to latent team skill parameters.
In their simplest form, these models treat matches as binary events, focusing only on the final outcome rather than the underlying scoreline.

In this work, we consider a family of Bradley–Terry–type models of increasing complexity, beginning with a basic binary formulation and progressively incorporating home effect and the possibility of draws.
These models serve both as a baseline for predictive performance and as a test for methodological validation, given their relatively low-dimensional parameter space and well-understood identifiability properties.

\subsection{General Formulation}
\label{sec:BT_general_formulation}

Before presenting the specific variants of the Bradley–Terry model, we introduce a general formulation that incorporates all its structural components: team skills, home effect, and the possibility of draws.
This serves as a conceptual starting point, from which different models naturally emerge through parameter restrictions and simplifying assumptions.

The Bradley-Terry model can be expressed in terms of team strengths and outcome probabilities.
For a match between team $i$ (home) and team $j$ (away), the strengths are defined as:
\begin{align}
    \label{eq:bt_general_home_strength}
    s_i &= \exp(\theta_i + \nu) \\
    \label{eq:bt_general_away_strength}
    s_j &= \exp(\theta_j),
\end{align}

\noindent where:
\begin{itemize}
    \item $\theta_i \in \mathbb{R}$ is the log-skill parameter for team $i$;
    \item $\nu \in \mathbb{R}$ is the home effect parameter;
    \item $s_i$ is the strength of team $i$ when playing at home;
    \item $s_j$ is the strength of team $j$ when playing away.
\end{itemize}

The probability of each outcome is then given by:
\begin{align}
    \label{eq:bt_general_prob_home}
    P(\text{team } i \text{ wins}) &= \frac{s_i}{s_i + s_j + s_{\text{tie}}} \\
    \label{eq:bt_general_prob_away}
    P(\text{team } j \text{ wins}) &= \frac{s_j}{s_i + s_j + s_{\text{tie}}} \\
    \label{eq:bt_general_prob_tie}
    P(\text{draw}) &= \frac{s_{\text{tie}}}{s_i + s_j + s_{\text{tie}}},
\end{align}

\noindent where $s_{\text{tie}} = \kappa \sqrt{s_i \cdot s_j}$ is the tie strength, with $\kappa \geq 0$ being the tie parameter, following the approach of \textcite{davidson_1970}.
When $\kappa = 0$, draws are not allowed and the model reduces to a binary outcome model, with the outcome $y_{i,j}$ following a Bernoulli distribution.
When $\kappa > 0$, the outcome $y_{i,j}$ follows a categorical distribution with these probabilities.

Given a set of matches $\mathcal{M}$, the likelihood can be computed as:
\begin{equation}
    \label{eq:bt_general_likelihood}
    \mathcal{L}(\boldsymbol{\theta}, \nu, \kappa ; \mathbf{y}) = \prod_{m \in \mathcal{M}} P(y_m \mid \theta_{i_m}, \theta_{j_m}, \nu, \kappa),
\end{equation}

\noindent where $y_m$ is the outcome of match $m$ between teams $i_m$ and $j_m$.

The Bradley-Terry model defined by equations \ref{eq:bt_general_prob_home}, \ref{eq:bt_general_prob_away}, and \ref{eq:bt_general_prob_tie} presents identifiability issues, with the likelihood remaining invariant under translations of all log-skill parameters by a constant $c$:
\begin{equation}
    \mathcal{L}(\boldsymbol{\theta}, \nu, \kappa ; \mathbf{y}) = \mathcal{L}(\boldsymbol{\theta} + c, \nu, \kappa ;\mathbf{y}),
\end{equation}

\noindent since:
\begin{align*}
    (\theta_i + c) + \nu - (\theta_j + c) &= \theta_i + \nu - \theta_j, \\
    \exp(\theta_i + c + \nu) &= \exp(c) \cdot \exp(\theta_i + \nu), \\
    \exp(\theta_j + c) &= \exp(c) \cdot \exp(\theta_j),
\end{align*}

\noindent and the constant $\exp(c)$ cancels out in the ratio of strengths.

To ensure model identifiability, we impose a sum-to-zero constraint on the log-skill parameters:
\begin{equation}
    \label{eq:sum_zero_cons_bt}
    \sum_{i=1}^{n} \theta_i = 0.
\end{equation}

This constraint is implemented by parametrizing the vector as:
\begin{equation}
    \boldsymbol{\theta} = \lrc{\theta_1, \theta_2, \ldots, \theta_{n-1}, -\sum_{i=1}^{n-1} \theta_i},
\end{equation}

\noindent where only the first $n-1$ parameters are estimated, and the last is determined by the constraint.

\subsection{Model Variants}
\label{sec:BT_model_variants}

Different models arise from simplifying assumptions about the parameters in the general formulation.
The models can be categorized based on two key dimensions:
\begin{enumerate}
    \item \textbf{Home Effect}: Present ($\nu \neq 0$) vs. Absent ($\nu = 0$);
    \item \textbf{Ties Allowed}: Present ($\kappa > 0$) vs. Absent ($\kappa = 0$).
\end{enumerate}

This section describes the four Bradley-Terry model variants (BT1--BT4) used to model match outcomes in football matches in the present work.
These models are organized from the simplest binary outcome model without home effect, BT1, to the most complex ternary outcome model with home effect, BT4.
The model names were chosen for simplicity, avoiding unnecessary repetition.

Model \textbf{BT1} is the simplest Bradley-Terry model, using a single log-skill parameter $\theta_i$ per team with no home effect and no ties allowed.
The log-odds of team $i$ winning against team $j$ is given by $\log(\text{odds}) = \theta_i - \theta_j$, which corresponds to a binary outcome model where $P(\text{team } i \text{ wins}) = \exp(\theta_i) / (\exp(\theta_i) + \exp(\theta_j))$.

Model \textbf{BT2} extends BT1 by incorporating a home effect parameter $\nu$.
The log-odds becomes $\log(\text{odds}) = \theta_i + \nu - \theta_j$ for the home team, capturing the additional effect of playing at home.
Like BT1, this model does not allow ties.

Model \textbf{BT3} extends BT1 by allowing ties through the introduction of a tie parameter $\kappa \geq 0$.
The tie strength is defined as $s_{\text{tie}} = \kappa \sqrt{s_i \cdot s_j}$, where $s_i = \exp(\theta_i)$ and $s_j = \exp(\theta_j)$.
This model does not include home effect, so the strengths are symmetric with respect to home and away contexts.

Model \textbf{BT4} combines the features of BT2 and BT3, incorporating both home effect and ties.
The home team's strength becomes $s_i = \exp(\theta_i + \nu)$, while the away team's strength remains $s_j = \exp(\theta_j)$.
The tie strength is $s_{\text{tie}} = \kappa \sqrt{s_i \cdot s_j}$, providing the most flexible specification that accounts for both home effect and the possibility of draws.

The taxonomy of these model variants is illustrated in Figure~\ref{fig:bt_models_taxonomy}.

\begin{figure}[ht]
    \centering
    \begin{tikzpicture}[
        level 1/.style = {sibling distance = 4cm, level distance = 2.5cm},
        level 2/.style = {sibling distance = 2cm, level distance = 2cm},
        every node/.style = {
            draw,
            rectangle,
            rounded corners,
            align=center,
            minimum width = 1.5cm,
            minimum height = 0.8cm,
            font = \scriptsize
        },
        root/.style = {
            minimum width = 3.5cm,
            minimum height = 1.5cm,
            font = \normalsize
        },
        leaf/.style = {
            fill = gray!20,
            font = \tiny
        }
    ]
        \node[root] {
            $\begin{aligned}
                s_i &= \exp(\theta_i + \nu) \\
                s_j &= \exp(\theta_j) \\
                s_{\text{tie}} &= \kappa \sqrt{s_i \cdot s_j}
            \end{aligned}$ \\
            General Formulation
        }
            child {
                node {
                    $\kappa = 0$ \\
                    No Ties
                }
                child {
                    node {
                        $\nu = 0$ \\
                        No Home \\
                        Effect
                    }
                    child {
                        node[leaf] {BT1}
                    }
                }
                child {
                    node {
                        $\nu \neq 0$ \\
                        With Home \\
                        Effect
                    }
                    child {
                        node[leaf] {BT2}
                    }
                }
            }
            child {
                node {
                    $\kappa > 0$ \\
                    Ties Allowed
                }
                child {
                    node {
                        $\nu = 0$ \\
                        No Home \\
                        Effect
                    }
                    child {
                        node[leaf] {BT3}
                    }
                }
                child {
                    node {
                        $\nu \neq 0$ \\
                        With Home \\
                        Effect
                    }
                    child {
                        node[leaf] {BT4}
                    }
                }
            };
    \end{tikzpicture}
    \caption{
        \textbf{Taxonomy of the four Bradley-Terry model variants}.
        The root node displays the complete specification with log-skill parameters ($\theta_i$), home effect ($\nu$), and tie parameter ($\kappa$).
        The first level of branching corresponds to whether ties are allowed ($\kappa > 0$) or not ($\kappa = 0$).
        The second level distinguishes models with ($\nu \neq 0$) and without ($\nu = 0$) home effect.
        Each leaf node represents a specific model variant, with the model identifier (BT1--BT4) indicating the final configuration after all simplifications have been applied.
    }
    \label{fig:bt_models_taxonomy}
\end{figure}


\section{Poisson model}
\label{sec:Poisson_model}

From a modelling perspective, Bradley-Terry models treat matches as binary or ternary events by collapsing the game score into a win/draw/loss indicator.
This abstraction makes it impossible for the model to distinguish between a dominant win and a narrow one, resulting in a loss of information about team performance.

In contrast, Poisson models \cite{maher_1982, dixon_1997, lee_1997, karlis_2000} offer an alternative by modelling the home and away goals directly.
This not only doubles the number of observed data per match, from an indicator to two count variables, but also allows the model to capture potential asymmetries between teams' performance.

\subsection{General Formulation}
\label{sec:Poisson_general_formulation}

The general formulation of the Poisson models presented in this section describes the most comprehensive structure considered in this work for modelling the number of goals scored by each team.
This parametrization includes distinct offensive and defensive effects, distinguishes between home and away contexts, and allows for the inclusion of a common effect term on both goal intensities, serving as the foundation for all model variants introduced later.

The general formulation for the log-intensity parameters can be expressed as:
\begin{align}
    \label{eq:poisson_general_home}
    \log(\lambda_{i,j,h}) &= \alpha_i + \beta_j + \nu + \eta \\
    \label{eq:poisson_general_away}
    \log(\lambda_{j,i,a}) &= \delta_j + \gamma_i + \eta,
\end{align}

where:
\begin{itemize}
    \item $\lambda_{i,j,h}$ is the expected number of goals scored by team $i$ (home) against team $j$ (away);
    \item $\lambda_{j,i,a}$ is the expected number of goals scored by team $j$ (away) against team $i$ (home);
    \item $\alpha_i$ is the offensive ability of team $i$ when playing at home;
    \item $\gamma_i$ is the defensive ability of team $i$ when playing at home;
    \item $\beta_j$ is the defensive ability of team $j$ when playing away;
    \item $\delta_j$ is the offensive ability of team $j$ when playing away;
    \item $\nu \in \mathbb{R}$ is the home effect parameter;
    \item $\eta \in \mathbb{R}$ is a common effect parameter that affects both teams equally.
\end{itemize}

The goals are then modelled as conditionally independent Poisson random variables:
\begin{align}
    \label{eq:poisson_general_dist_home}
    g_{i,j} \mid \lambda_{i,j} &\sim \text{Poisson}(\lambda_{i,j}) \\
    \label{eq:poisson_general_dist_away}
    g_{j,i} \mid \lambda_{j,i} &\sim \text{Poisson}(\lambda_{j,i}),
\end{align}

\noindent where $g_{i,j}$ and $g_{j,i}$ are the goals scored by teams $i$ and $j$ respectively.

Given a set of matches $\mathcal{M}$, with home goals $G_h$ and away goals $G_a$, the likelihood can be computed as:
\begin{equation}
    \label{eq:poisson_general_likelihood}
    \mathcal{L}(\boldsymbol{\theta} ; G_h, G_a) = \prod_{m \in \mathcal{M}} \frac{\lambda_{m_h}^{g_{h_m}} \cdot \exp(-\lambda_{m_h})}{g_{h_m}!} \cdot \frac{\lambda_{m_a}^{g_{a_m}} \cdot \exp(-\lambda_{m_a})}{g_{a_m}!},
\end{equation}

\noindent where $g_{h_m}$ and $g_{a_m}$ are the total goals of home and away teams in match $m$, respectively, and $\lambda_{m_h}$ and $\lambda_{m_a}$ are the expected number of goals for the home and away teams in match $m$.

The general formulation presented in equations \ref{eq:poisson_general_home} and \ref{eq:poisson_general_away} presents identifiability issues.
Specifically, the likelihood remains invariant under certain transformations of the parameters.

When $\boldsymbol{\alpha}$, $\boldsymbol{\beta}$, $\boldsymbol{\gamma}$, and $\boldsymbol{\delta}$ are all independent, adding a constant $c$ to all attack parameters and subtracting the same constant from the defence parameters, while keeping $\nu$ and $\eta$ unchanged, does not change the likelihood:
\begin{equation}
    \mathcal{L}(\boldsymbol{\alpha}, \boldsymbol{\beta}, \boldsymbol{\gamma}, \boldsymbol{\delta}, \nu, \eta ; G_h, G_a) = \mathcal{L}(\boldsymbol{\alpha} + c, \boldsymbol{\beta} - c, \boldsymbol{\gamma} - c, \boldsymbol{\delta} + c, \nu, \eta ; G_h, G_a).
\end{equation}

This occurs because:
\begin{align*}
    (\alpha_i + c) + (\beta_j - c) + \nu + \eta &= \alpha_i + \beta_j + \nu + \eta \\
    (\delta_j + c) + (\gamma_i - c) + \eta &= \delta_j + \gamma_i + \eta.
\end{align*}

The constant $c$ cancels out in the additions and subtractions, leaving the log-intensity parameters unchanged, and thus the likelihood remains invariant.
Similarly, adding a constant $c$ on $\eta$ and subtracting $c$ on both $\alpha$ and $\delta$ or both $\beta$ and $\gamma$ keep the likelihood unchanged:
\begin{align*}
    \mathcal{L}(\boldsymbol{\alpha}, \boldsymbol{\beta}, \boldsymbol{\gamma}, \boldsymbol{\delta}, \nu, \eta ; G_h, G_a) &  = \mathcal{L}(\boldsymbol{\alpha} - c, \boldsymbol{\beta}, \boldsymbol{\gamma}, \boldsymbol{\delta} - c, \nu, \eta + c ; G_h, G_a) \\
    & = \mathcal{L}(\boldsymbol{\alpha}, \boldsymbol{\beta} - c, \boldsymbol{\gamma} - c, \boldsymbol{\delta}, \nu, \eta + c ; G_h, G_a).
\end{align*}
To ensure model identifiability, we impose sum-to-zero constraints parametrizing the vectors as:
\begin{align}
    \label{eq:sum_to_zero_alpha}
    \boldsymbol{\alpha} &= \lrc{\alpha_1, \alpha_2, \ldots, \alpha_{n-1}, -\sum_{i=1}^{n-1} \alpha_i} \\
    \label{eq:sum_to_zero_beta}
    \boldsymbol{\beta} &= \lrc{\beta_1, \beta_2, \ldots, \beta_{n-1}, -\sum_{i=1}^{n-1} \beta_i} \\
    \label{eq:sum_to_zero_delta}
    \boldsymbol{\delta} &= \lrc{\delta_1, \delta_2, \ldots, \delta_{n-1}, -\sum_{i=1}^{n-1} \delta_i} \\
    \label{eq:sum_to_zero_gamma}
    \boldsymbol{\gamma} &= \lrc{\gamma_1, \gamma_2, \ldots, \gamma_{n-1}, -\sum_{i=1}^{n-1} \gamma_i}
\end{align}

\noindent where only the first $n-1$ parameters are estimated, and the last is determined by the constraint.

In models where $\eta$ is set equal to zero, the constraints given by Equations \ref{eq:sum_to_zero_beta} and \ref{eq:sum_to_zero_delta} can be removed, since the restrictions of Equations \ref{eq:sum_to_zero_alpha} and \ref{eq:sum_to_zero_gamma} are already sufficient to ensure identifiability.
However, in our implementation, we maintain all constraints regardless of whether $\eta$ is present.
This approach ensures consistency across all model variants and enhances interpretability by fixing a reference point (zero mean) for all parameter vectors, allowing parameters to be interpreted as deviations from the average team ability.

\subsection{Model Variants}
\label{sec:Poisson_model_variants}

Different models arise from simplifying assumptions about the parameters in equations \ref{eq:poisson_general_home} and \ref{eq:poisson_general_away}. The models can be categorized based on three key dimensions:
\begin{enumerate}
    \item \textbf{Offence-Defence Parametrization}: Single parameter ($\boldsymbol{\alpha} = \boldsymbol{\delta}$ and $\boldsymbol{\beta} = \boldsymbol{\gamma}$, with $\boldsymbol{\beta} = -\boldsymbol{\alpha}$) vs. Two parameters ($\boldsymbol{\alpha}$ and $\boldsymbol{\beta}$ independent, with $\boldsymbol{\alpha} = \boldsymbol{\delta}$ and $\boldsymbol{\beta} = \boldsymbol{\gamma}$) vs. Four parameters ($\boldsymbol{\alpha}$, $\boldsymbol{\beta}$, $\boldsymbol{\gamma}$, and $\boldsymbol{\delta}$ all independent);
    \item \textbf{Home Effect}: Present ($\nu \neq 0$) vs. Absent ($\nu = 0$);
    \item \textbf{Common Effect}: Present ($\eta \neq 0$) vs. Absent ($\eta = 0$).
\end{enumerate}

This section describes all the Poisson model variants (P1--P10) used to model goal scoring in football matches in the present work.
These models are organized from the most general formulation, ranging from the simplest single-parameter model, P1, to the most complex four-parameter model with home effect and common effect, P10.
The model names were chosen for simplicity, avoiding unnecessary repetition.

Model \textbf{P1} uses a single skill parameter $\boldsymbol{\alpha}$ per team, with $\log(\lambda_{i,j}) = \alpha_i - \alpha_j$ and $\log(\lambda_{j,i}) = \alpha_j - \alpha_i = -\log(\lambda_{i,j})$.
It assumes identical offensive and defensive abilities, with no home effect or common effect.

Model \textbf{P2} extends P1 by adding a home effect parameter $\nu$, giving $\log(\lambda_{i,j}) = \alpha_i + \nu - \alpha_j$ for the home team and $\log(\lambda_{j,i}) = \alpha_j - \alpha_i$ for the away team.

Model \textbf{P3} introduces separate offensive ($\boldsymbol{\alpha}$) and defensive ($\boldsymbol{\beta}$) parameters, with $\log(\lambda_{i,j}) = \alpha_i + \beta_j$.
It assumes identical abilities at home and away, with no home effect or common effect.

Model \textbf{P4} combines P3's two-parameter structure with home effect: $\log(\lambda_{i,j}) = \alpha_i + \beta_j + \nu$ for the home team and $\log(\lambda_{j,i}) = \alpha_j + \beta_i$ for the away team.

Model \textbf{P5} uses four independent parameter sets: $\boldsymbol{\alpha}$, $\boldsymbol{\gamma}$ for home abilities and $\boldsymbol{\beta}$, $\boldsymbol{\delta}$ for away abilities, with $\log(\lambda_{i,j,h}) = \alpha_i + \beta_j$ and $\log(\lambda_{j,i,a}) = \gamma_i + \delta_j$.
This parametrization already captures home effect implicitly through the difference between home and away parameter sets, making an explicit home effect parameter $\nu$ redundant and potentially causing identifiability issues.
Model P5 includes no common effect.

Models \textbf{P6} through \textbf{P10} extend P1 through P5, respectively, by adding a common effect parameter $\eta$ to both log-intensities: $\log(\lambda_{i,j}) = \log(\lambda_{i,j}^{\text{base}}) + \eta$ and $\log(\lambda_{j,i}) = \log(\lambda_{j,i}^{\text{base}}) + \eta$, where the base structure follows the corresponding model P1--P5.
This formulation is inspired by the bivariate Poisson distribution proposed by \textcite{holgate_1964}, where shared components appear in both rate parameters.
In our models, the common effect parameter $\eta$ acts as a shared additive term on both log-intensities, uniformly scaling both teams' scoring rates by a constant factor $\exp(\eta)$ across all matches.
This parameter captures systematic deviations in the overall level of scoring that are not explained by team strengths alone, such as league-wide tactical trends or seasonal effects, rather than match-specific factors such as weather conditions or pitch quality, which would require a match-varying parameter.

The taxonomy of these model variants is illustrated in Figure~\ref{fig:poisson_models_taxonomy}.

\begin{figure}[ht]
    \centering
    \begin{tikzpicture}[
        level 1/.style = {sibling distance = 6cm, level distance = 2.5cm},
        level 2/.style = {sibling distance = 3cm, level distance = 2cm},
        level 3/.style = {sibling distance = 1.5cm, level distance = 1.8cm},
        every node/.style = {
            draw,
            rectangle,
            rounded corners,
            align=center,
            minimum width = 1.25cm,
            minimum height = 0.8cm,
            font = \scriptsize
        },
        root/.style = {
            minimum width = 3.5cm,
            minimum height = 1.5cm,
            font = \normalsize
        },
        leaf/.style = {
            fill = gray!20,
            font = \tiny
        }
    ]
        \node[root] {$\begin{aligned}
            \log(\lambda_{i,j,h}) &= \alpha_i + \beta_j + \nu + \eta \\
            \log(\lambda_{j,i,a}) &= \delta_j + \gamma_i + \eta
        \end{aligned}$\\General Formulation}
            child {
                node {
                    $\boldsymbol{\beta} = -\boldsymbol{\alpha}$ \\
                    $\boldsymbol{\delta} = \boldsymbol{\alpha}$ \\
                    $\boldsymbol{\gamma} = \boldsymbol{\beta} = -\boldsymbol{\alpha}$ \\
                    Single Parameter
                }
                child {
                    node {
                        $\nu = 0$ \\
                        No Home \\
                        Effect
                    }
                    child {
                        node[leaf] {
                            $\eta = 0$ \\ P1
                        }
                    }
                    child {
                        node[leaf] {
                            $\eta \neq 0$ \\ P6
                        }
                    }
                }
                child {
                    node {
                        $\nu \neq 0$ \\
                        With Home \\
                        Effect
                    }
                    child {
                        node[leaf] {
                            $\eta = 0$ \\ P2
                        }
                    }
                    child {
                        node[leaf] {
                            $\eta \neq 0$ \\ P7
                        }
                    }
                }
            }
            child {
                node {
                    $\boldsymbol{\alpha} = \boldsymbol{\delta}$ \\
                    $\boldsymbol{\beta} = \boldsymbol{\gamma}$ \\
                    Two Parameters
                }
                child {
                    node {
                        $\nu = 0$ \\
                        No Home \\
                        Effect
                    }
                    child {
                        node[leaf] {
                            $\eta = 0$ \\ P3
                        }
                    }
                    child {
                        node[leaf] {
                            $\eta \neq 0$ \\ P8
                        }
                    }
                }
                child {
                    node {
                        $\nu \neq 0$ \\
                        With Home \\
                        Effect
                    }
                    child {
                        node[leaf] {
                            $\eta = 0$ \\ P4
                        }
                    }
                    child {
                        node[leaf] {
                            $\eta \neq 0$ \\ P9
                        }
                    }
                }
            }
            child {
                node {
                    All independent \\
                    Four Parameters
                }
                child {
                    node {
                        $\nu = 0$ \\
                        No Home \\
                        Effect
                    }
                    child {
                        node[leaf] {
                            $\eta = 0$ \\ P5
                        }
                    }
                    child {
                        node[leaf] {
                            $\eta \neq 0$ \\ P10
                        }
                    }
                }
            };
    \end{tikzpicture}
    \caption{
        \textbf{Taxonomy of the ten Poisson model variants}.
        The root node displays the complete four-parameter specification with independent offensive and defensive parameters for home and away contexts ($\boldsymbol{\alpha}$, $\boldsymbol{\beta}$, $\boldsymbol{\gamma}$, $\boldsymbol{\delta}$), home effect ($\nu$), and common effect ($\eta$).
        The first level of branching corresponds to the parametrization structure: single parameter (where offensive and defensive abilities are constrained to be equal), two parameters (where home and away abilities are constrained to be equal), or four parameters (where all abilities are independent).
        The second level distinguishes models with ($\nu \neq 0$) and without ($\nu = 0$) home effect.
        The third level separates models with ($\eta \neq 0$) and without ($\eta = 0$) common effect.
        Each leaf node represents a specific model variant, with the model identifier (P1--P10) indicating the final configuration after all simplifications have been applied.
    }
    \label{fig:poisson_models_taxonomy}
\end{figure}


\section{Adaptation to Player-Level Models}
\label{sec:player_level_adaptation}

The models described in the previous sections operate at the club level, where team abilities are modelled directly as team-specific parameters.
However, these models can be adapted to operate at the player level.
This adaptation provides a more granular understanding of team performance by decomposing it into individual player contributions.

Instead of estimating team abilities directly, we estimate individual player abilities and aggregate them, weighting by playing time, to obtain team-level strengths.
For a given match, let $\mathcal{P}_h$ and $\mathcal{P}_a$ denote the sets of players participating in the home and away teams, respectively, and let $m_p$ denote the number of minutes played by player $p$ in that match.
The team-level skill is given by:
\begin{equation}
    \label{eq:team_skill_aggregation}
    s_{\text{team}} = \sum_{p \in \mathcal{P}} \exp(\theta_p) \cdot \frac{m_p}{90},
\end{equation}

\noindent where $\theta_p$ is the log-skill parameter for player $p$, and $\mathcal{P}$ is either $\mathcal{P}_h$ or $\mathcal{P}_a$ depending on the team.
The normalization by 90 minutes ensures that a player's contribution is proportional to their playing time relative to the standard match duration.

When a player is sent off (red card), the team plays with fewer players for the remainder of the match.
To handle this situation, we introduce a synthetic player with an estimable ability parameter $\theta_{\text{synthetic}}$ that represents the impact of playing with fewer players.
Specifically, if a player is sent off at minute $t_{\text{off}}$, we add a synthetic player with $m_{\text{synthetic}} = 90 - t_{\text{off}}$ minutes.
If two or more players receive a red card, the same treatment is applied, increasing the value of $m_{\text{synthetic}}$.
The synthetic player's ability $\theta_{\text{synthetic}}$ is estimated by the model and quantifies the impact of the red card on team performance.
This approach allows the model to learn how much a sending-off affects team strength.

As club-level models, player-level models also require identifiability constraints.
The same sum-to-zero constraint is imposed on the player skill parameters:
\begin{equation}
    \label{eq:player_sum_zero_constraint}
    \sum_{p=1}^{n} \theta_p = 0,
\end{equation}

\noindent where $n$ is the total number of players, including the synthetic.
This constraint ensures that the model is identifiable and that player abilities are measured relative to the average player ability.


\section{Inference Procedure}
\label{sec:inference_procedure}

This section describes the prior distributions, data inputs, and computational settings used for posterior inference across all models.
All models are fitted within a Bayesian framework using Stan \cite{carpenter_2017} via the CmdStanPy interface \cite{cmdstanpy}, with the No-U-Turn Sampler (NUTS) as the Hamiltonian Monte Carlo algorithm.

All club-level models use independent $\N(0, 1)$ priors for every parameter.
This applies to team log-skill parameters ($\boldsymbol{\theta}$, $\boldsymbol{\alpha}$, $\boldsymbol{\beta}$, $\boldsymbol{\gamma}$, $\boldsymbol{\delta}$), the home effect ($\nu$), the tie parameter ($\kappa$), and the common effect parameter ($\eta$), whenever each is present in the model.
For player-level Poisson models, the prior on individual skill parameters is set to $\N(0, 0.1)$.
Since team strength is obtained by aggregating the contributions of 11--16 players (Equation~\ref{eq:team_skill_aggregation}), the tighter per-player prior ensures numerical stability during posterior sampling.
Global parameters ($\nu$ and $\eta$) retain the $\N(0, 1)$ prior at the player level, as their dimensionality does not change.
Table~\ref{tab:priors} summarises the prior distributions used across all models.

\begin{table}[htbp]
    \centering
    \begin{tabular}{llcc}
        \toprule
        Parameter & Description & Club-Level & Player-Level \\
        \midrule
        $\boldsymbol{\theta}$, $\boldsymbol{\alpha}$, $\boldsymbol{\beta}$, $\boldsymbol{\gamma}$, $\boldsymbol{\delta}$ & Skill parameters & $\N(0, 1)$ & $\N(0, 0.1)$ \\
        $\nu$ & Home effect & $\N(0, 1)$ & $\N(0, 1)$ \\
        $\kappa$ & Tie parameter (BT) & $\N(0, 1)$ & --- \\
        $\eta$ & Common effect & $\N(0, 1)$ & $\N(0, 1)$ \\
        \bottomrule
    \end{tabular}
    \caption{
        \textbf{Prior distributions for all model parameters}.
        Skill parameters receive a tighter prior at the player level to ensure numerical stability when aggregating individual contributions into team strength.
        The tie parameter $\kappa$ is only present in Bradley-Terry models, which are not extended to the player level.
    }
    \label{tab:priors}
\end{table}

Club-level models receive, for each match, the identifiers of the home and away teams and the number of goals scored by each.
For Bradley-Terry models, the goals are collapsed into a categorical outcome (home win, draw, or away win).
No information about individual players, substitutions, or playing time is used at this level.
Player-level models require a richer data structure.
For each match, the model receives the list of players participating in each team, identified by their unique CBF ID, along with the number of minutes played by each player and the total goals scored by each team.
This information is derived from the temporal segmentation of match reports described in Section~\ref{sec:data_treatment}, which tracks substitutions and red cards to determine each player's exact playing time.
When a player receives a red card, the synthetic player described in Section~\ref{sec:player_level_adaptation} is assigned the remaining minutes of the match.

The scope of data used for fitting differs between the two analysis contexts.
For club-level model evaluation, the rolling-window scheme described in Section~\ref{sec:evaluation_methodology} fits the model incrementally: at each evaluation point, the model receives all \textit{Brasileirão Série A} matches played up to that point in the current season, starting from 50 matches (approximately 5 complete match-days) and progressing through the season in increments of 10 matches.
For player-level analysis, the model is fitted on the complete set of matches from the selected period: all 380 matches of the 2025 season for the single-season analysis and all 2,280 matches from 2020 to 2025 for the multi-season analysis.
In both cases, only \textit{Brasileirão Série A} data is used.

Club-level models are fitted with 4 parallel chains, each running 1,000 warmup iterations and 1,000 sampling iterations, yielding 4,000 posterior samples per fit.
The sampler uses an adaptation target acceptance probability of $\delta = 0.95$ and a maximum tree depth of 10.
This configuration is identical to the one used for SBC validation (Section~\ref{sec:sbc_experimental_design}).
Player-level models require a more intensive sampling configuration due to the substantially higher dimensionality of the parameter space.
Each fit uses 4 parallel chains with 2,500 warmup iterations and 7,500 sampling iterations per chain, yielding 30,000 posterior samples.
The extended warmup allows the sampler to adequately adapt its mass matrix to the high-dimensional geometry, while the larger number of sampling iterations ensures sufficient effective sample size across hundreds of parameters.


\section{Evaluation Methodology}
\label{sec:evaluation_methodology}

This section presents the framework used to evaluate the predictive performance of all club-level models on real \textit{Brasileirão Série A} data from 2020 to 2025.

We evaluate model predictions using a rolling-window prediction scheme.
After each match-day $g$, the model is fitted using all matches from rounds 1 to $g$, and predictions are generated for all future matches.
For match-level metrics, we evaluate the predicted probabilities for round $g+1$; for season-level metrics, the model simulates the remaining $380 - g$ matches to produce predictive distributions of accumulated points for each team at every future match-day.
This process begins at round 50 (approximately 5 complete match-days) to ensure sufficient training data, and continues through round 370, yielding predictions for approximately 330 matches per season evaluated out-of-sample.
All reported metrics are averaged across evaluation points within each season and then across the five seasons.

Predictions are generated via Monte Carlo integration over the posterior distribution.
For each future match, 1,000 parameter vectors are sampled from the posterior.
For Bradley-Terry models, match outcomes are sampled directly from the categorical distribution implied by the model.
For Poisson models, home and away goals are simulated independently from Poisson distributions with rates determined by the sampled parameters, and match outcome probabilities are computed as the empirical frequencies of home wins, draws, and away wins.
Point totals are accumulated across simulated matches, yielding predictive distributions from which quantile-based intervals are extracted.

To contextualize model performance, we compare against two baseline predictors:
\begin{itemize}
    \item \textbf{Naive 1}: Assigns uniform probabilities $(1/3, 1/3, 1/3)$ to all outcomes, representing a completely uninformative forecast.
    \item \textbf{Naive 2}: Assigns probabilities based on historical outcome frequencies observed in the current season, reflecting the marginal distribution of results up to match-day $g$.
\end{itemize}

We assess predictive quality using four complementary scoring rules, all negatively oriented (lower values indicate better predictions).
The first three metrics evaluate match-level outcome predictions, while the fourth assesses season-level point forecasts.
Let $n$ denote the number of evaluated matches.
For each match $t$, let $\mathbf{z}_t = (z_{H,t}, z_{D,t}, z_{A,t})$ be the one-hot encoded outcome vector, where $H$, $D$, and $A$ denote home victory, draw, and away victory, respectively.
Let $\mathbf{p}_t = (p_{H,t}, p_{D,t}, p_{A,t})$ be the corresponding predicted probabilities.

The \textbf{Brier Score} (BS) measures the accuracy of probabilistic predictions for categorical outcomes, computed as the mean squared error between predicted probabilities and observed outcomes:
\begin{equation}
    \text{BS} = \frac{1}{n} \sum_{t=1}^{n} \sum_{j \in \{H, D, A\}} (z_{j,t} - p_{j,t})^2.
    \label{eq:brier_score}
\end{equation}

The \textbf{Log Score} (LS), also known as the negative log-likelihood, is a strictly proper scoring rule that heavily penalizes confident but incorrect predictions:
\begin{equation}
    \text{LS} = -\dfrac{1}{n} \sum_{t=1}^{n} \sum_{j \in \{H, D, A\}} z_{j,t} \log(p_{j,t}).
    \label{eq:log_score}
\end{equation}

The \textbf{Ranked Probability Score} (RPS) accounts for the ordinal nature of match outcomes (home win $\succ$ draw $\succ$ away win), penalizing predictions that are ``further'' from the true outcome:
\begin{equation}
    \text{RPS} = \dfrac{1}{2n} \sum_{t=1}^{n} \sum_{k=1}^{2} \left( \sum_{j=1}^{k} (z_{j,t} - p_{j,t}) \right)^2.
    \label{eq:rps}
\end{equation}

Finally, the \textbf{Interval Score} (IS) evaluates the quality of predictive intervals for cumulative points.
Unlike the previous metrics that assess single-match predictions, IS measures how well the model forecasts the range of possible point totals for each team at future match-days.

A good predictive interval should be both \emph{narrow} (providing precise information) and \emph{calibrated} (containing the true value at the stated confidence level).
The IS captures this trade-off: for a confidence level $(1-\alpha)$, let $[l, u]$ be the predicted interval and $y$ the observed points:
\begin{equation}
    \text{IS}_\alpha = (u - l) + \frac{2}{\alpha}(l - y)\mathbf{1}_{y < l} + \frac{2}{\alpha}(y - u)\mathbf{1}_{y > u}.
    \label{eq:interval_score}
\end{equation}
The first term penalizes wide intervals, while the second and third terms penalize coverage failures, which are cases where the observed value falls outside the interval.
The factor $2/\alpha$ ensures that narrower intervals (smaller $\alpha$) incur larger penalties when they fail to cover the true value.

We aggregate IS across all 20 teams, all future match-days, and 19 symmetric prediction intervals ranging from 5\% to 95\% coverage.

For season-level analysis of specific \textit{Brasileirão} seasons, we apply a complementary approach focused on a single well-performing model.
After each match-day, the model is refitted using all available match results, and 10,000 Monte Carlo simulations are used to forecast the remaining matches.
Each iteration samples parameters from the posterior distribution and simulates all unplayed matches according to the model's likelihood, producing a complete final standings table.

The primary output is a full probability matrix assigning each team a probability for each final position (1--20).
From this matrix, we can derive quantities such as:
\begin{itemize}
    \item Probability of winning the championship;
    \item Probability of qualifying for \textit{Copa Libertadores da América} (\textit{Libertadores}), South America's premier continental club competition;
    \item Probability of qualifying for \textit{Copa Sul-Americana} (\textit{Sul-Americana}), South America's secondary continental club competition;
    \item Probability of relegation.
\end{itemize}

Additionally, posterior distributions of team strength parameters are obtained after each match-day, providing interpretable summaries of relative team quality.

